% III. Проектиране на системата.
\section{Проектиране на системата}
\label{sec:design}

Системата е разделена на четири слоя, така че всеки от тях да може да бъде проверен
самостоятелно. Слоят за придобиване на данни доставя байтове, независимо дали са от
жив интерфейс или от записан файл. Слоят за разчитане превръща байтовете в структури
на протокола. Слоят за проверка взема решение за валидността на сертификатната верига.
Слоят за допускане решава дали една нова връзка да бъде обслужена. Общата схема на
слоевете и потокът от данни между тях се дават на \TODO{фигура с блоковата схема на
системата, с легенда по А.9}.

\subsection{Слой за придобиване и разчитане}
\label{subsec:parser}

Придобиването е скрито зад един интерфейс с два реализатора: единият чете от жив
интерфейс през libpcap, другият от записан файл. Тестовете използват втория, защото
той е детерминиран.

Разчитането следва структурата на записите и на съобщенията от \cite{rfc8446}.
Съществено проектно решение тук е границата на видимостта. В TLS 1.3 само първата
част от ръкостискането е в явен вид, а останалата е шифрована с ключове, които
пасивният наблюдател няма. Затова анализаторът докладва това, което наистина се
наблюдава, и обозначава изрично като неизвестно всичко останало, вместо да го
извежда по подразбиране. Пълният списък на наблюдаемите и на скритите параметри се
дава в \TODO{таблица с наблюдаеми и скрити параметри на ръкостискането}.

Разчитането работи изцяло върху невладеещи изгледи към получения буфер. Нито една
структура не притежава байтовете, които описва. Това премахва копирането по пътя на
анализа и същевременно поставя ясно изискване към времето на живот на буфера, което
се проверява от тестовете.

\subsection{Слой за проверка на сертификатната верига}
\label{subsec:validator}

Проверката се извършва на два етапа. Първият изгражда пътя от сертификата на края до
доверен корен, като разглежда всички възможни продължения, а не само първото
намерено. Вторият проверява за всеки сертификат по пътя ограниченията, изисквани от
\cite{rfc5280}: срок на валидност, подпис на издателя, основни ограничения,
предназначение на ключа, ограничения върху имената и съответствие на името на края.

Проверката на отмяната е отделена от проверката на пътя, защото има различна цена и
различен режим на отказ. Поддържат се два източника: списък на отменените сертификати
и протоколът за състояние по \cite{rfc6960}. Политиката при недостъпен източник е
проектно решение с последици за сигурността и се задава явно, а не се подразбира.
Разгледаните варианти и приетият се описват в \TODO{вариантите на политиката при
недостъпен източник на информация за отмяна}.

\subsection{Слой за допускане на връзки}
\label{subsec:admission_design}

Слоят за допускане стои пред сървъра и решава дали да отдели ресурс за нова връзка.
Проектиран е като верига от независими механизми, всеки от които може да бъде включен
или изключен, без това да засегне останалите. Такава организация е необходима, защото
измерването в Раздел~\ref{sec:method} изисква всеки механизъм да бъде оценен и
поотделно, и в съчетание с другите.

Отчитането на таблицата с връзките е общо за цялата верига. То поддържа броя на
незавършените и на установените връзки, разпределението им по източници и времето,
прекарано във всяко състояние. Тези величини са едновременно вход за решенията и
изход за измерването.

\subsection{Модел на заплахата и граници на системата}
\label{subsec:threat_model}

Системата приема, че наблюдателят няма достъп до частните ключове и не може да
дешифрира защитената част от обмена. Приема също, че опитите с лавинен трафик се
провеждат в среда, която е изцяло под контрола на автора. Извън модела на заплахата
остават атаките срещу самата машина, върху която работи системата, както и атаките
срещу приложния слой над транспорта. Пълното изложение на модела и на приетите
допускания се дава в \TODO{таблица с модела на заплахата: актьор, способност, мярка}.
